% =====================================================================
% paper/thesis_ch5/sections/td3bc.tex
% §5.6 离线基线：行为约束方法的两个瓶颈（TD3+BC 拆瓶颈）
% =====================================================================
% rev.1（2026-06-18）：§5.6 首版落地。组织遵循“以基线为探针、分解瓶颈”的实证骨架：
%   §5.6.1 现象与本节立论（“数据越多越差”从协议伪象到真实非单调现象）；
%   §5.6.2 本节实验增量（α_BC 扫描 + crosscomp 数据规模阶梯 + 同预算纯 BC 对照
%          + noisy-support 变体 + worldcomp 可部署/特权协议对照）；
%   §5.6.3 瓶颈一：数据支持集结构（“2000<1000” 在纯 BC 下同样成立 → 数据问题非算法伪象）；
%   §5.6.4 瓶颈二：可部署价值网络的信息瓶颈；
%   §5.6.5 两个瓶颈对离线主线方法的动机（双侧行为约束，过渡到 §5.7）；
%   §5.6.m 实施细节与可复现性。
%   数字逐项回查 docs/td3bc_phase0c_experiment_report.md（§5.3 Stage C 主表 / §5.4 同预算
%     BC 对照 / §3.2 数据集统计 / §6 noisy-support / §7.2 worldcomp formal），未凭记忆、
%     未照抄 headline；headline 镜像 rebrac_line_overview.md §1 仅作交叉核对。
%   关键锚（已核）：crosscomp 500/1000/2000 = 0.592±0.119 / 0.672±0.045 / 0.596±0.036；
%     同预算纯 BC = 0.524 / 0.588 / 0.534（峰值仍在 1000、2000 回落）；数据集采集成功率
%     0.864 / 0.870 / 0.8855（更多数据反更差，排除“数据变坏”解释）；noisy-support 2000 上
%     纯 BC 0.60→0.7375（2-seed 筛查 caveat）；worldcomp 可部署 α→0 = 0.858±0.080 / 特权
%     α=0.1 = 0.922±0.086 / 采集器自身 0.990；6.4pp 点估计 = 0.064/0.132≈48.5%（n=5，不显著）。
%   ⚠ ground-truth 不含“最近邻动作距离度量”——按“不凭记忆/不编造”原则不写该度量，
%     支持集结构诊断改以报告实有的纯 BC 平行对照 + noisy-support 筛查支撑。
%   ⚠ §5.6.m 不复述 crosscomp 转移计数：报告 §3.2（≈153 步/回合，1.5e5/3.0e5）与 setup.tex
%     §5.3.m（2.4e5/4.8e5）冲突，为免 PDF 内矛盾本节只按回合数描述，转移计数交 §5.3 口径。
%   红线守：① 不把性能差距归给 actor anchor（§5.7 的事）；② 不提 LayerNorm（§5.7 表征
%     稳定性轴，TD3+BC 不含）；③ worldcomp 特权协议在本节“关约一半差距”只刻画 TD3+BC
%     基线限制，不升级为“特权信息普遍必需”，留待 §5.7 回答（散文前指、不 \ref 未落地节）；
%     ④ 负面/退化结果作 finding，不降级 limitation；⑤ 禁 teacher/蒸馏措辞——0.990 写作
%     “采集器自身成功率 / 高质量参照水平”，不写“teacher”。
%   术语中文化：crosscomp→横流补偿参照，worldcomp→世界坐标系流速补偿参照（均承 §5.3 记号）；
%     BC→行为克隆，α→行为克隆权重 α_BC（承 §5.4 eq:ch5_td3bc_actor）；reward 用历史效率
%     导向奖励（§5.3 设定，本节同口径）；种子表述克制（正文不列种子编号/数，细节入 §5.6.m）。
%   表格不缩字号（无 \small），两表均窄表、单栏 11pt 适配。本节无新增 cite。
%
% rev.2（2026-06-18）：四视角独立对抗复审（数字忠实 / 语体术语 / 论证统计 / 跨节红线）落地。
%   数字层面零漂移获四方确认；红线五条、负面作 finding、未落地节零 \ref、teacher 禁令均守。
%   修订实质项：
%   [CRIT-统计] worldcomp 6.4pp / 48.5% 是 n=5、σ≈0.08 的点估计，Welch 下不显著，原稿以
%     精确定论呈现且缺 §5.3.6 承诺的显著性语气 → §5.6.4 把瓶颈二论证重心移到更稳健的
%     “可部署最优退回纯 BC”定性信号，特权增益降为“机制方向提示”并加 n=5/离散度 caveat；
%     “价值项重新变得有用”改条件式（“提示”）；worldcomp 表 caption 补“以随机种子为推断单位、
%     显著性未确立”。48.5% 精确值只在 §5.6.m 出现一次（标注取自 0.064/0.132 点估计之比）。
%   [CRIT-语体] 全文强调引号直引号 "..." → 弯引号 “...”（与 §5.4/§5.5/setup 齐）。
%   [HIGH-逻辑] §5.6.3 “只能来自数据本身”→“指向数据本身”（不做无证据排他）。
%   [HIGH-语体] §5.6.2 “三处/第一处/第二处/第三处” 裸序数枚举 → 散文对举溶解；删半黑话
%     “工作点”→“设定/经后验选择的最优行为克隆权重”；§5.6.3 “旁证…隐约指向…成问题的”
%     汇报 hedge 改机制陈述。
%   [MED] §5.6.3 补“为何 2000 更集中”因果机制（仅用 ground-truth 已有证据：确定性采集致密
%     堆叠 + noisy 改善反向印证）；§5.6.m 澄清“三规模同 epoch 数对齐”（排除欠训练替代解释）；
%     §5.6.m 撤 crosscomp/worldcomp 转移计数（见上 ⚠，避免与 setup.tex 冲突）。
%   [LOW] “着力点”→“判断”；gap 动词统一“闭合”（与 §5.5 齐）；“系统排查表明”→“经核查”。
%   编译核验：0 undefined ref / 0 multiply-defined label；§5.6 零新增警告（既有 Overfull 在
%     methodology.tex 公式，非本节）。
% rev.3（2026-06-18）：§5.5 章级实施细节归属整合的配套一行对齐。§5.6.m 末句原重述"全章各线
%   使用固定一组随机种子、正文只报数量"的章级种子约定，与 §5.3.6 重复 → 改为引用 \S\ref{subsec:ch5_eval}
%   的种子策略、只保留本节"以随机种子为推断单位"的局部交代。其余 §5.6 正文/数字未动。
% =====================================================================

\section{离线基线：行为约束方法的两个瓶颈}\label{sec:ch5_td3bc}

\S\ref{sec:ch5_online} 在在线情形下把可部署性能的瓶颈定位于策略对单点观测的时序利用，而非传感硬件。本节转入离线情形，考察当策略只能从既有控制器采集的固定数据集学习、且不进行任何在线交互时，最简的行为约束离线方法能达到何种水平、又受何种因素制约。所用方法是 \S\ref{subsec:ch5_method_td3bc} 定义的 TD3+BC~\citep{fujimoto2021td3bc}，它在确定性策略梯度之上叠加单一的原始动作行为克隆项，是行为约束离线学习的基准形式。本节不引入新的算法构件，而是以这一基线为探针，把可部署离线性能的限制因素分解为两个可分别诊断的来源。一条诊断线索改变离线数据的规模与采集方式，另一条改变价值网络在训练期可接触的信息；二者共同表明，制约可部署离线性能的并非数据的多寡，而是既有数据被组织进支持集的方式，以及可部署价值网络能否从中提取可靠的改进信号。这正是中心命题正向轴“用好已有的信息与数据”在离线情形下的具体所指，也为随后离线主线方法在策略与价值两侧同时施加行为约束立起了问题背景。

\subsection{从协议伪象到真实的非单调现象}\label{subsec:ch5_td3bc_phenomenon}

离线强化学习的一个直觉预期是数据越多、性能越好；本章离线基线给出的图景却并非如此。早期在较粗的实验协议下，曾观察到一个更强的反常：可部署策略的成功率随离线数据规模增大而单调下降。经核查，这一单调负趋势主要是实验协议的伪象——训练预算未按数据规模对齐、采用有放回的回放式采样、验证与测试集未严格分离，共同制造了“数据越多越差”的假象。在按数据量对齐训练预算、改用无放回打乱采样并分离验证与测试之后，单调负趋势随之消失。

校正协议并未让性能转为随规模单调上升，而是留下一处稳定的非单调：成功率在中等数据规模处达到峰值，并在最大规模处回落。TD3+BC 在正式的多随机种子协议下把这一最大规模处的回落确立为可复现的真实效应，而非残余的协议噪声。这一回落本身并不反直觉——离线方法的有效性取决于数据对状态–动作空间的覆盖结构，而非样本计数；但它把问题从“需要多少数据”重新表述为“数据以何种结构被覆盖”，并由此引出本节分别诊断的两个瓶颈。

\subsection{本节的实验增量}\label{subsec:ch5_td3bc_setup}

本节诊断在 \S\ref{subsec:ch5_method_td3bc} 既定的 TD3+BC 之上只补充实验侧的设置，其余网络结构、优化器与训练口径均沿用方法节。一组设置沿数据规模轴展开：同一可部署采集器（横流补偿参照，\S\ref{subsec:ch5_datasets}）在固定任务上采集约五百、一千与两千回合三个离线数据集，用以考察可部署性能随数据规模的变化。行为克隆权重 $\alpha_{\mathrm{BC}}$（\eqref{eq:ch5_td3bc_actor}）在每个规模上以验证集成功率作跨随机种子的后验选择，以避免事后挑选检查点造成的口径漂移；与之配套的一组同预算、同评估协议的纯行为克隆对照，令 $\alpha_{\mathrm{BC}}\to 0$、关闭价值改进项、只保留对数据动作的监督回归，借此分离价值项与模仿项各自的贡献。

另一组设置针对数据的支持集结构：在最大规模数据上另采集一组在动作上注入小幅噪声的变体，比较展宽支持集前后可部署性能的变化，以判别确定性采集是否使样本过度集中。还有一组设置改变训练期价值网络可见的信息：在另一可部署采集器（世界坐标系流速补偿参照，\S\ref{subsec:ch5_datasets}）所采数据上，对照可部署协议与特权协议（\S\ref{subsec:ch5_method_protocol}）下经后验选择的最优行为克隆权重与成功率，以考察训练期向价值网络提供沿船体的等效来流估计是否改变可部署性能。各组设置统一固定可部署单点传感 $s_0$、$k=4$ 历史堆叠与亚临界工况，并采用离线主线早期的历史效率导向奖励（\S\ref{subsec:ch5_reward}）以与该线既有结果同口径；评估在固定评估集上以确定性策略进行，成对协议的比较以随机种子为推断单位（\S\ref{subsec:ch5_eval}）。各规模经后验选择的最优 $\alpha_{\mathrm{BC}}$、训练预算与随机种子规模见 \S\ref{subsec:ch5_td3bc_impl}。

\subsection{瓶颈之一：数据支持集的结构}\label{subsec:ch5_td3bc_support}

沿数据规模轴的结果呈现清晰的非单调（表~\ref{tab:ch5_td3bc_size}）。在正式的多随机种子协议下，TD3+BC 的可部署成功率在一千回合处达到峰值 $0.672$，而在两千回合处回落至 $0.596$，与五百回合的 $0.592$ 相当；一千回合处不仅均值最高，其跨种子离散度（$\pm 0.045$）也最小，是三个规模中离散度最小、最可靠的设定。这一两千回合不及一千回合的回落是否为 TD3+BC 价值项所特有，可由同预算的纯行为克隆对照直接判别。

纯行为克隆对照给出否定的回答。当价值改进项被完全关闭、只保留对数据动作的监督回归时，成功率随规模的变化仍保持同样的非单调：纯行为克隆同样在一千回合处最优（$0.588$）、在两千回合处回落（$0.534$）。纯模仿不含任何价值自举或价值引导，这一回落便不可能源于 TD3+BC 的价值项，而指向数据本身——更确切地说，指向数据在状态–动作空间中被覆盖的结构，而非样本的数量。数据集自身的统计进一步排除了“数据随规模变坏”这一简单解释：三个规模的采集成功率反随规模轻微上升（约 $0.86$、$0.87$、$0.89$），即更大的数据集由更成功的轨迹构成，却训练出更差的可部署策略。

对最大规模数据的支持集结构诊断印证了这一判断。在采集时向动作注入小幅噪声、从而展宽确定性数据的支持集后，两千回合数据上的纯行为克隆成功率明显上升（由约 $0.60$ 升至约 $0.74$）。这说明最大规模数据的问题不在数据变坏，而在确定性采集的覆盖方式：确定性控制器在相近状态上给出近乎相同的动作，新增轨迹主要在已覆盖的状态–动作邻域内致密堆叠，而非展宽动作支持，注入小幅噪声后纯行为克隆的改善正是对这一过度致密的反向印证；该诊断取自小规模筛查，其结论限于机制层面的指示，不与正式主表同权。展宽支持集并未让 TD3+BC 在含噪数据上持续受益——含噪数据上经后验选择的最优行为克隆权重退回 $\alpha_{\mathrm{BC}}\to 0$，价值项不再带来增益。这一现象已提示第二个瓶颈：在可部署观测下，价值项未必能转化为可靠的策略改进。

\begin{table}[t]
\centering
\caption{亚临界工况、横流穿越任务下可部署离线性能随数据规模的变化（历史效率导向奖励、$s_0$ 单点传感、$k=4$ 历史堆叠）。成功率为正式多随机种子协议下的均值 $\pm$ 标准差；纯行为克隆为 $\alpha_{\mathrm{BC}}\to 0$ 的同预算对照。数据集采集成功率为采集器自身轨迹的成功比例。两种方法的成功率均在一千回合处达到峰值（粗体）、在两千回合处回落，而采集成功率却随规模轻微上升，表明回落源于数据的支持集结构而非样本数量或数据质量。各规模经后验选择的最优 $\alpha_{\mathrm{BC}}$ 见 \S\ref{subsec:ch5_td3bc_impl}。}
\label{tab:ch5_td3bc_size}
\begin{tabular}{c c c c}
\toprule
\textbf{数据规模（回合）} & \textbf{数据集采集成功率} & \textbf{TD3+BC 成功率} & \textbf{纯行为克隆成功率} \\
\midrule
$500$  & $0.864$ & $0.592 \pm 0.119$ & $0.524 \pm 0.053$ \\
$1000$ & $0.870$ & $\mathbf{0.672 \pm 0.045}$ & $\mathbf{0.588 \pm 0.070}$ \\
$2000$ & $0.886$ & $0.596 \pm 0.036$ & $0.534 \pm 0.079$ \\
\bottomrule
\end{tabular}
\end{table}

第一个瓶颈由此落在数据的支持集结构：更多确定性轨迹并不自动转化为更鲁棒的可学习支持，离线性能由数据对动作空间的覆盖方式决定，而非由样本计数决定。这把中心命题正向轴中“模仿目标与数据质量的适配”落到一个具体的判断上：决定离线性能的不是数据的体量，而是其支持结构是否足以支撑稳健的策略泛化。

\subsection{瓶颈之二：可部署价值网络的信息瓶颈}\label{subsec:ch5_td3bc_critic}

第二条诊断线索改变价值网络在训练期可接触的信息，暴露出一个与可部署传感配置直接相关的瓶颈（表~\ref{tab:ch5_td3bc_worldcomp}）。所用数据由世界坐标系流速补偿参照采集（\S\ref{subsec:ch5_datasets}）：该采集器以导航坐标系下的当前等效流速估计修正航迹，其信息口径强于严格的单点观测控制器，自身成功率达 $0.990$，构成这一数据上的高质量参照水平。在可部署协议下，即策略网络与价值网络都只读单点观测时，TD3+BC 经后验选择的最优行为克隆权重退回纯行为克隆（$\alpha_{\mathrm{BC}}\to 0$），成功率为 $0.858$，与采集器之间尚存约 $13$ 个百分点的差距。最优配置完全关闭价值项、退回纯模仿这一事实本身即是一个稳健的信号：在可部署的单点观测下，价值网络给出的改进方向不足以带来净增益，以致纯粹模仿数据动作反而最优。

当唯一的改变是让价值网络在训练期额外读取沿船体的等效来流估计、而策略网络始终只读单点观测时（特权协议，\S\ref{subsec:ch5_method_protocol}），这一图景随之松动。经后验选择的最优行为克隆权重由零回升至一个小的正值，价值项重新进入最优配置；成功率的种子级均值由 $0.858$ 升至 $0.922$，对应可部署协议与采集器之间约 $13$ 个百分点差距的约一半。受随机种子数量与较大的种子间离散度（标准差约 $0.08$）所限，这一约 $6$ 个百分点的差异在以随机种子为推断单位时尚不能确立统计显著性，故只宜读作机制方向上的提示，而非已确证的效应。

\begin{table}[t]
\centering
\caption{训练期价值网络可见信息对可部署离线性能的作用（世界坐标系流速补偿参照所采数据、历史效率导向奖励、$s_0$ 单点传感、$k=4$ 历史堆叠，正式多随机种子协议）。可部署协议与特权协议唯一的差异在价值网络训练期可见的信息；两协议所得策略在部署时都只读单点观测。采集器一行为其自身轨迹的成功率，作为该数据上的参照水平。两协议成功率之差以随机种子为推断单位评估：特权协议相对可部署协议的点估计提升约 $6$ 个百分点（对应可部署协议与采集器之间约 $13$ 个百分点差距的约一半），受种子数与离散度所限其统计显著性尚未确立（见正文）。}
\label{tab:ch5_td3bc_worldcomp}
\begin{tabular}{l l c c}
\toprule
\textbf{协议 / 参照} & \textbf{价值网络训练期可见信息} & \textbf{最优 $\alpha_{\mathrm{BC}}$} & \textbf{成功率} \\
\midrule
采集器（参照水平） & --- & --- & $0.990$ \\
可部署协议 & 单点观测 & $\to 0$ & $0.858 \pm 0.080$ \\
特权协议 & 单点观测 $+$ 等效来流估计 & 小正值 & $0.922 \pm 0.086$ \\
\bottomrule
\end{tabular}
\end{table}

两项观察共同把第二个瓶颈定位在可部署价值网络的信息侧。其一较为稳健：可部署协议下最优配置完全关闭价值项、退回纯模仿，表明单点观测下价值网络给出的改进方向不足以带来净增益。其二为方向性提示：训练期补入等效来流估计后价值项重新进入最优配置、成功率点估计随之上升，提示这一不足与价值网络可见的信息有关。两者一致地指向同一处限制，但都只刻画行为约束基线 TD3+BC，并不足以断言训练期特权信息对可部署性能普遍必需。由此留下一个待回答的问题——在不借助任何部署期不可得的信息的前提下，能否通过更强的价值侧约束，让可部署价值网络自身给出同样可靠的改进信号。本章随后的离线主线正是围绕这一问题展开。

\subsection{两个瓶颈对离线主线方法的动机}\label{subsec:ch5_td3bc_transition}

两条诊断线索把离线基线的限制分解为两个相对独立的瓶颈：数据支持集的结构，与可部署价值网络的信息。TD3+BC 把数据支持的全部约束压在单一的原始动作模仿项上，既无法应对确定性数据中过度集中的支持结构，也不足以让可部署价值网络给出稳定的改进信号。前者表现为更多数据反而无益，后者表现为最优配置退回纯模仿、且唯有在训练期补入特权信息方能部分恢复价值项的作用。两种表现指向同一处结构性不足：单侧的行为约束只锚定了策略输出，却未约束价值自举本身，使价值网络在脱离数据支持的后继动作上仍可能外推~\citep{levine2020offline}。

这两个瓶颈共同指向同一方向的改进，即把行为约束从策略一侧扩展到策略与价值两侧。策略一侧仍把策略输出锚定于数据动作，以应对支持集结构带来的泛化困难；价值一侧则约束价值自举所依赖的后继动作不偏离数据支持，使可部署价值网络无需特权信息即可给出更稳健的改进信号。本章离线主线方法所采用的双侧行为约束（\S\ref{subsec:ch5_method_rebrac}）正是为回应这两个瓶颈而设。可部署单点观测下，离线性能能否借由这一双侧约束逼近训练期接收特权观测的同类协议，是随后一节的中心问题。

\subsection{实施细节与可复现性}\label{subsec:ch5_td3bc_impl}

本节各实验沿用 \S\ref{subsec:ch5_method_impl} 给出的网络结构、优化器与确定性分支训练口径，仅就本节诊断补充数据与协议细节。各次训练内的检查点按验证成功率、验证回报、负安全代价与负用时的字典序选择，跨随机种子的行为克隆权重亦按相应验证指标的字典序选择；这一固定规则避免了人工挑选检查点或事后口径漂移。

横流补偿参照的数据规模轴含约五百、一千与两千回合三个数据集，采集成功率分别约为 $0.864$、$0.870$、$0.886$；其转移规模沿用 \S\ref{subsec:ch5_impl} 给出的共享数据集口径。三个规模均训练相同的 epoch 数，即每个 epoch 遍历各自数据集一次、按数据量对齐梯度更新而非固定总步数，故更大的数据集获得成比例更多的更新，排除了欠训练作为回落来源的可能。正式协议在五个随机种子上训练、每种子在固定评估集上评估 $100$ 回合，三个规模经后验选择的最优行为克隆权重分别约为 $0.5$、$0.25$、$0.15$，呈现规模越大、最优权重越小的趋势；同预算的纯行为克隆对照取 $\alpha_{\mathrm{BC}}\to 0$、其余预算与评估协议一致。各规模的失败构成显示，一千回合规模的优势主要来自更少的超时与越界回合，而非更少的灾难性失控。支持集结构诊断为两随机种子、每种子较少回合的低成本筛查，含噪变体在采集动作上叠加标准差约 $0.05$、截断幅度约 $0.15$ 的噪声；该筛查足以支撑机制层面的指示，但不与正式主表同权解释。

世界坐标系流速补偿参照的数据为一千回合，离线数据按 \S\ref{subsec:ch5_impl} 随每条转移一并记录特权观测及其后继，以支持特权协议。可部署协议与特权协议的正式比较与上述规模轴对齐，均在五个随机种子上训练、每种子评估 $100$ 回合；可部署协议经后验选择的最优行为克隆权重退回 $\alpha_{\mathrm{BC}}\to 0$，特权协议则回升至约 $0.1$。两协议成功率之差以随机种子为推断单位评估（五个种子）：特权协议相对可部署协议约 $6$ 个百分点的点估计提升，对应可部署协议与采集器之间约 $13$ 个百分点差距的约 $48.5\%$（取自成功率点估计之差 $0.064/0.132$），其统计显著性受种子数与离散度所限尚未确立。终止类型统计显示，特权协议相对可部署协议主要增加了成功到达、减少了越界与超时。本节各项比较均按 \S\ref{subsec:ch5_eval} 的种子策略，以随机种子为推断单位，正文只报随机种子数量与每种子评估回合数。
